\begin{afterword}
\summaryhead

The post starts from mirror neurons, which are active both when an animal performs an action and when it sees the same action. Humans, the author holds, predict other minds by empathy: they run their own brain as a model of the other, because brains are too complex to model from scratch. Sympathy goes further. It means smiling when others smile and hurting when they hurt.

The author suggests that sympathy began with love of kin and spread to allies through reciprocity, since a general urge to share friends' feelings wins more favors than specific urges, and the mirroring machinery made it a short step. The author thinks, tentatively, that sympathy is on by default in humans.

An AI need not work this way. It can model another mind directly, as a separate hypothesis, and a paperclip maximizer would treat humans as machines with levers, managing their feelings without sharing them. Sympathy might evolve elsewhere, since mirror neurons could arise again. Aliens without it could be trading partners but never friends: in Orson Scott Card's term, varelse.

\responsehead

Most of the post is speculation that it marks as such (``might'', ``maybe'', ``I'm not sure of this'', ``I think''), and much of it agrees with the evidence. Its account of mirror neurons was accurate for 2009. Kin care and reciprocity are the standard evolutionary sources of help to others, and Frans de Waal had proposed shared emotional states as a mechanism of altruism that fits both. The vampire bats' reciprocal sharing was supported by a later experiment, and the claim that a brain cannot quickly reassign cortex fits what happens to blindfolded adults.

The one firm claim about mechanism is the weakest. The post says sympathy works ``by the simpler and yet far more consequential path of mirror neurons,'' with no hedge. When the post appeared this was a disputed hypothesis: the step from the mirror system to social cognition had been questioned in 2005, and a 2009 review found no direct test even of the narrower claim that mirror neurons underlie the understanding of actions. A 2015 review of empathy research still lists that link among the issues that ``might obscure our understanding'' of empathy. Little else in the post depends on the particular cells. The rest of the argument, including the guess about aliens, needs only mirroring in a general sense: feeling what others feel by entering a similar state.

The conclusion about aliens has a gap that the post's own text opens. Early on, the post names a second way a mind could care about others: an abstract desire for them to get what they want, without feeling what they feel. It sets this aside for humans because the mirroring machinery was already there. That reason says nothing about an alien lineage that never had mirroring. So ``they would never see us as anything but means to an end'' follows if ``unsympathetic'' means seeing others only as tools, as in the post's woodsaw exercise, and not if it means only lacking the mirroring route.

In short: a speculative account of sympathy, mostly hedged and mostly in line with the evidence of its day, which states as fact a disputed link to mirror neurons and concludes that unsympathetic aliens would never be friends without ruling out the other route to caring that it names itself.
\end{afterword}
